\chapter{骑宠、攻防混合与最终伤害公式}

本章讨论战斗系统中另一个容易被经验说法简化过度的部分：骑宠状态下，人物与宠物的攻击、防御、敏捷究竟如何混合；而这种混合又如何继续进入物理伤害、属性修正与会心附加伤害。若只看面板，很容易误以为骑宠只是“骑上去加一点能力”；但从源码角度看，骑宠实际上改变了多个核心变量的来源层级，因此它不仅影响数值大小，也改变了数值进入公式的路径。

\section{骑宠攻击力与防御力混合}

在当前研究口径下，骑宠三围的读取主要对应 \code{BATTLE_adjustRidePet3A()}，并以前文已说明的 \code{_BATTLE_NEWPOWER} 分支为主。这里首先要区分是否装备远程武器，因为远程和近战的混合方式并不相同。

设人物攻击力、防御力分别为 $A_c,D_c$，骑宠攻击力、防御力分别为 $A_p,D_p$。若攻击方未骑宠，则物理攻击直接取
\begin{equation}
    A = A_c.
\end{equation}
若攻击方已骑宠，则当装备远程武器时，攻击力混合为
\begin{equation}
    A = A_c + 0.4A_p,
    \label{eq:ride-atk-ranged}
\end{equation}
而当装备近战武器或空手时，攻击力混合为
\begin{equation}
    A = 0.8A_c + 0.8A_p.
    \label{eq:ride-atk-melee}
\end{equation}
式~\eqref{eq:ride-atk-ranged} 与式~\eqref{eq:ride-atk-melee} 的差异说明，骑宠近战并不是单纯“由宠物代打”，而是人物与宠物攻击共同占据较高权重；相反，骑宠远程则明显偏向人物本体，只把宠物作为额外附加项。

防御侧的结构更特殊。若防守方未骑宠，则在 \code{BATTLE_DamageCalc()} 中会先取
\begin{equation}
    D = 0.70D_c.
    \label{eq:def-no-ride}
\end{equation}
若防守方已骑宠，则首先在骑宠混合函数中得到
\begin{equation}
    D_{ride}=0.7D_c+0.7D_p,
\end{equation}
随后再进入伤害函数乘以一次 $0.70$，于是最终用于物伤计算的防御量为
\begin{equation}
    D = 0.70D_{ride}=0.49D_c+0.49D_p.
    \label{eq:def-ride}
\end{equation}
因此，在当前编译口径下，骑宠防御既不是人物防御与宠物防御的简单平均，也不是只取宠物，而是对两者做同权线性组合后再统一衰减一次。

\section{骑宠敏捷混合}

骑宠敏捷的处理与攻击、防御不同，它必须区分“攻击侧读取”与“防守侧读取”，还必须区分远程武器与近战武器。设人物战斗速度敏为 $Q_c$，骑宠战斗速度敏为 $Q_p$。则在攻击侧，若使用远程武器，读取为
\begin{equation}
    Q_{atk}^{(ranged)} = 0.8Q_c + 0.2Q_p,
    \label{eq:ride-quick-ranged}
\end{equation}
若使用近战武器或空手，则读取为
\begin{equation}
    Q_{atk}^{(melee)} = 0.2Q_c + 0.8Q_p.
    \label{eq:ride-quick-melee}
\end{equation}
而在防守侧，读取进一步偏向宠物，写成
\begin{equation}
    Q_{def} = 0.1Q_c + 0.9Q_p.
    \label{eq:ride-quick-def}
\end{equation}
这说明骑宠并不仅仅改变“谁出手”，还改变“速度由谁主导”。从直觉上说，骑宠近战明显更偏宠物，骑宠远程明显更偏人物，而骑宠防守几乎可以看作宠物敏捷主导。正因为如此，同一只宠物在“骑近战”和“骑远程”两种配置下，既可能表现出完全不同的先手性质，也可能在后续伤害、命中与反击中体现出不同的角色定位。

\section{物理伤害三段式公式}

设前两节得到的攻击量与防御量分别记为 $A$ 与 $D$。在当前 \code{_BATTLE_NEWPOWER} 分支下，物理伤害的基础体可分成三段。若
\begin{equation}
    D > A,
\end{equation}
则基础物伤只在极低范围内波动，可整理为
\begin{equation}
    damage = \operatorname{RAND}(0,1).
    \label{eq:damage-stage1}
\end{equation}
若
\begin{equation}
    D \le A < \frac{8}{7}D,
\end{equation}
则进入第二段：
\begin{equation}
    damage = \operatorname{RAND}\left(0,\frac{A}{16}\right).
    \label{eq:damage-stage2}
\end{equation}
若
\begin{equation}
    A \ge \frac{8}{7}D,
\end{equation}
则进入第三段，记
\begin{equation}
    K_0 = \operatorname{RAND}\left(0,\frac{A}{8}\right)-\frac{A}{16},
\end{equation}
则基础物伤为
\begin{equation}
    damage = 2(A-D)+K_0.
    \label{eq:damage-stage3}
\end{equation}
其中当前源码常量对应
\begin{equation}
    DAMAGE\_RATE=2.0, \qquad D\_8=\frac{1}{8}, \qquad D\_{16}=\frac{1}{16}.
\end{equation}
从结构上看，式~\eqref{eq:damage-stage1}--\eqref{eq:damage-stage3} 构成了一个典型的分段物伤模型：当攻击显著低于防御时，伤害几乎被压到 $0$ 或 $1$；当攻击刚刚接近防御时，仍停留在小幅随机区间；只有当攻击高于 $8D/7$ 之后，伤害才开始表现为对 $(A-D)$ 的线性函数并叠加轻微随机扰动。

\section{防御动作、全防与偏转的介入顺序}

除了前文的防御力工作值以外，当前物理伤害流程里还存在一条独立的“防御链”。它不直接改写 \code{CHAR_WORKDEFENCEPOWER}，而是在基础伤害已经算出之后，再按动作状态或专门判定去压低最终伤害。

其中最基础的一层是\emph{防御动作}。若防守方当前动作是 \code{BATTLE_COM_GUARD}，且未处于混乱，则源码不会简单按固定比例减伤，而是调用 \code{BATTLE_GuardAdjust()} 做一次六档随机压缩：
\begin{equation}
    damage_{guard} \in \{0,0.1,0.2,0.3,0.4,0.5\}\cdot damage,
\end{equation}
其概率分布分别为
\begin{equation}
    (25\%,25\%,20\%,15\%,10\%,5\%).
\end{equation}
也就是说，防御并不是“固定减半”，而是以完全挡下到保留一半伤害之间的离散随机分布实现。若压缩后伤害降到 $0$，系统还会继续把这次结果标记为 \code{BATTLE_RET_ALLGUARD}，在表现层显示为防御成功。

第二层是\emph{偏转 / 招架}。若编译路径启用了 \code{_EQUIT_ARRANGE}，则在普通防御动作之外，还会额外调用 \code{BATTLE_ArrangeCheck()}。一旦成立，本次伤害会被直接压到原值的 $10\%$，并返回独立的偏转结果码。对玩家阅读而言，这一层更接近“招架成功后大幅卸力”，而不是普通防御动作本身。

第三层是\emph{忽视防御}。若编译路径启用了 \code{_EQUIT_NEGLECTGUARD}，且攻击方装备或状态给出了 \code{CHAR_WORKNEGLECTGUARD}，则它并不是去跳过“防御动作”，而是在基础物伤计算前先按比例削减防守方防御量。若忽视防御比例记为 $r\%$，则进入物伤公式的防御量会先变成
\begin{equation}
    D' = \left(1-\frac{r}{100}\right)D.
\end{equation}
因此，忽视防御影响的是三段式物伤的分段位置，而不是后面的防御动作随机压缩。

把这几层放到同一条顺序里，可以更准确地写成：先由攻击力与防御力工作值得到基础物伤，再经过属性修正与结界减伤，然后再检查防御动作、偏转和伤害反应等后置环节。于是，玩家口中的“被格挡了”在源码里其实可能对应三种完全不同的机制：普通防御动作的随机减伤、偏转 / 招架的 $10\%$ 压缩，或吸收 / 反射 / 消失等伤害反应。它们在日志表现上都可能让单次伤害显著下降，但进入公式的位置并不相同。


\section{属性克制、场地倍率与结界减伤}

基础物伤只是整个伤害流程的第一层。对当前实现而言，式~\eqref{eq:damage-stage1}--\eqref{eq:damage-stage3} 得到的结果还要继续经过属性克制、场地倍率与结界减伤，然后才进入额外扰动与后续表现层。与很多玩家经验不同，这一段并不是“看主属性相克后乘一个固定倍率”，而是一个由\emph{五维属性向量、双线性相克矩阵与场地比值}共同组成的计算链。

\subsection{攻击方与防守方的五维属性向量}

在当前 \code{_BATTLE_NEWPOWER} 口径下，物理伤害使用的属性值由 \code{BATTLE\_GetAttr()} 直接从单位自身的四属性工作值读取，而不会因为骑宠状态再把人物与宠物属性做一次平均。也就是说，骑宠会混合攻击、防御与敏捷，但\textbf{不会在当前主分支里混合物理属性向量}。若把攻击方四属性写成
\begin{equation}
    (p_E,p_W,p_F,p_{Wi}),
\end{equation}
把防守方四属性写成
\begin{equation}
    (q_E,q_W,q_F,q_{Wi}),
\end{equation}
则源码会先把负值裁到 $0$，再补出无属性分量
\begin{equation}
    p_N = \max\{100-p_E-p_W-p_F-p_{Wi},0\},
    \qquad
    q_N = \max\{100-q_E-q_W-q_F-q_{Wi},0\}.
    \label{eq:none-attr}
\end{equation}
于是攻守双方真正进入属性克制计算的，分别是五维向量
\begin{equation}
    \bm{p}=(p_E,p_W,p_F,p_{Wi},p_N),
    \qquad
    \bm{q}=(q_E,q_W,q_F,q_{Wi},q_N).
\end{equation}

式~\eqref{eq:none-attr} 说明，只要四属性之和没有达到 $100$，剩余部分就会自动落入无属性；若四属性和已经超过 $100$，则无属性被截成 $0$，而不会出现负的无属性值。

\subsection{属性克制的源码矩阵}

记式~\eqref{eq:damage-stage1}--\eqref{eq:damage-stage3} 的基础物伤为 $damage_0$。源码会先把攻击方五维属性按比例拆到这份基础物伤上，再与防守方五维属性逐项相乘，并乘上三种常量
\begin{equation}
    AJ_{\mathrm{up}}=1.5,
    \qquad
    AJ_{\mathrm{same}}=1.0,
    \qquad
    AJ_{\mathrm{down}}=0.6.
\end{equation}
若按 $(E,W,F,Wi,N)$ 的顺序排列攻守属性，则源码等价于使用如下相克矩阵
\begin{equation}
    M=
    \begin{bmatrix}
        1.0 & 1.5 & 1.0 & 0.6 & 1.5 \\
        0.6 & 1.0 & 1.5 & 1.0 & 1.5 \\
        1.0 & 0.6 & 1.0 & 1.5 & 1.5 \\
        1.5 & 1.0 & 0.6 & 1.0 & 1.5 \\
        0.6 & 0.6 & 0.6 & 0.6 & 1.0
    \end{bmatrix}.
    \label{eq:attr-matrix}
\end{equation}
于是属性克制后的物伤主体可以直接写成
\begin{equation}
    damage_{\mathrm{attr}}
    = damage_0\,\frac{\bm{p}^{\mathsf T}M\bm{q}}{10000}.
    \label{eq:damage-attr-core}
\end{equation}

式~\eqref{eq:attr-matrix} 给出的口径有几个实战上很容易记反的点。其一，土克水、水克火、火克风、风克土；相克链并不是很多玩家口口相传的另一种写法。其二，同属性与“隔一位”的属性关系在当前物伤克制里都按 $1.0$ 处理。其三，\emph{无属性并不是中立常数}：元素打无属性按 $1.5$ 算，无属性打任意元素按 $0.6$ 算，只有无属性打无属性才是 $1.0$。

若只看纯属性对纯属性，则式~\eqref{eq:damage-attr-core} 可以压缩成表~\ref{tab:physical-attr-pure}。

\begin{table}[H]
    \centering
    \caption{物理伤害中纯属性对纯属性的源码倍率}
    \label{tab:physical-attr-pure}
    \begin{tabularx}{\textwidth}{>{\raggedright\arraybackslash}p{1.8cm} *{5}{>{\centering\arraybackslash}X}}
        \toprule
        攻 \textbackslash{} 守 & 土 & 水 & 火 & 风 & 无 \\
        \midrule
        土 & $1.0$ & $1.5$ & $1.0$ & $0.6$ & $1.5$ \\
        水 & $0.6$ & $1.0$ & $1.5$ & $1.0$ & $1.5$ \\
        火 & $1.0$ & $0.6$ & $1.0$ & $1.5$ & $1.5$ \\
        风 & $1.5$ & $1.0$ & $0.6$ & $1.0$ & $1.5$ \\
        无 & $0.6$ & $0.6$ & $0.6$ & $0.6$ & $1.0$ \\
        \bottomrule
    \end{tabularx}
\end{table}

\subsection{场地倍率不是加法，而是攻守比值}

属性克制并不是这一层的终点。源码接着会读取战场属性 \code{field\_att} 与场地强度 \code{att\_pow}，分别对攻击方与防守方算一次场地倍率，然后取二者比值乘回伤害。设当前场地属性为 $X\in\{E,W,F,Wi\}$，场地强度为 $s$，则对任意一方属性向量 $\bm{u}$，其场地倍率可写成
\begin{equation}
    \phi_X(\bm{u}) = 0.5 + 0.5\,u_X\,\frac{s}{10000},
    \label{eq:field-ratio}
\end{equation}
若当前没有属性场地，则源码回到默认值
\begin{equation}
    \phi_{\mathrm{none}}(\bm{u}) = 0.5.
\end{equation}
于是属性与场地共同作用后的结果是
\begin{equation}
    damage_1 = damage_{\mathrm{attr}}\,\frac{\phi_X(\bm{p})}{\phi_X(\bm{q})}.
    \label{eq:damage-after-field}
\end{equation}

式~\eqref{eq:damage-after-field} 的一个直接推论是：场地不是简单给攻击方加成，也不是简单给防守方减伤，而是对攻守双方\emph{同时各算一次倍率再取比值}。因此，同一个火场里，火高的攻击方确实会放大输出，但火高的防守方也会通过分母把这部分收益抵消掉一部分。

\subsection{四属性结界的减伤公式}

在 \code{_PROFESSION_ADDSKILL} 分支下，物伤经过式~\eqref{eq:damage-after-field} 后，还会再检查防守方是否挂有四属性结界。这里的实现有三个值得单独指出的源码细节。

第一，结界检查是按
\begin{equation}
    E \rightarrow W \rightarrow F \rightarrow Wi
\end{equation}
的 \code{if-else if} 顺序进行，因此同一时刻即使理论上挂了多个结界，也只会命中第一条成立的分支。第二，减伤读取的是\emph{攻击方原始固定属性工作值}，而不是式~\eqref{eq:damage-after-field} 里已经参与过场地和相克计算的属性向量。第三，结界高位里存放的“强度值”在当前物伤减伤公式中并不直接进入比例，只起到“该结界是否开启”的开关作用。

若记命中的结界属性为 $k\in\{E,W,F,Wi\}$，记攻击方对应的固定属性工作值为 $r_k$，则当前源码下的结界减伤可写成
\begin{equation}
    damage_2 = damage_1\left(1-\frac{r_k}{200}\right).
    \label{eq:boundary-reduce}
\end{equation}
若没有任何结界成立，则直接有
\begin{equation}
    damage_2 = damage_1.
\end{equation}

式~\eqref{eq:boundary-reduce} 说明，结界并不是“固定减半”或“按结界强度减伤”，而是按攻击方同属性点数做线性压缩。例如攻击方火属性固定值为 $40$，而防守方当前触发的是火结界，则物伤会再乘上 $1-40/200=0.8$。

\subsection{属性反转如何进入本章伤害公式}

第~\ref{chap:dynamic-attr} 章已经说明，属性反转并不是直接改写某个“永久主属性”，而是通过标志位控制当前固定属性快照的交换。在伤害层面，这意味着它影响的不是某个后置乘区，而是本节前面几条公式的\emph{输入量本身}。

若某单位当前有效固定属性由式~\eqref{eq:reverse-map} 从
\begin{equation}
    (E,W,F,Wi) \longmapsto (F,Wi,E,W)
\end{equation}
交换而来，那么在当前主分支下至少有三处会随之改变。

第一，\code{BATTLE\_GetAttr()} 读取到的攻击方或防守方属性向量 $\bm{p},\bm{q}$ 会改成反转后的版本，因此式~\eqref{eq:damage-attr-core} 中的双线性相克值会整体改变。第二，式~\eqref{eq:field-ratio} 中用于计算场地倍率的对应属性分量 $u_X$ 也会改成反转后的版本，因此同一块场地上的攻守比值可能随反转而跳变。第三，式~\eqref{eq:boundary-reduce} 中的 $r_k$ 直接读取攻击方当前固定属性工作值，所以结界减伤也会跟着反转后的属性一起变化。

因此，若把“反转前”的伤害记为
\begin{equation}
    damage_2 = \mathcal{D}(\bm{p},\bm{q},r_k),
\end{equation}
把“反转后”的伤害记为
\begin{equation}
    damage_2^{(rev)} = \mathcal{D}(\bm{p}^{(rev)},\bm{q}^{(rev)},r_k^{(rev)}),
\end{equation}
则差异不止发生在某一步，而是会同时穿过
\begin{equation}
    \bm{p}^{\mathsf T}M\bm{q},
    \qquad
    \frac{\phi_X(\bm{p})}{\phi_X(\bm{q})},
    \qquad
    1-\frac{r_k}{200}
\end{equation}
这三层结构。

这里还要补一个只从第十章公式本身不容易看出的时序细节。当前物伤函数并不重新检查“反转标志此刻是开是关”，它只读取\emph{当前这份已经算好的固定属性工作值}。因此，若出现第~\ref{chap:dynamic-attr} 章讨论过的过渡情形，即“反转标志已经被清掉，但新的固定属性重算尚未发生”，那么本章中的 $\bm{p},\bm{q},r_k$ 仍会继续使用那份尚未恢复的反转态快照。这也是一些实战现象看上去像“反转效果多保留了一段时间”的根源。

\subsection{两个完整算例}

为了避免把式~\eqref{eq:damage-attr-core}--\eqref{eq:boundary-reduce} 看成抽象符号，这里给出两个按当前源码口径直接展开的例子。

\paragraph{例 1：属性克制与场地比值共同放大伤害}
设基础物伤 $damage_0=120$，攻击方属性向量为 $\bm{p}=(60,20,10,0,10)$，防守方属性向量为 $\bm{q}=(10,50,20,0,20)$，当前为土场地，场地强度 $s=80$。按式~\eqref{eq:damage-attr-core} 与式~\eqref{eq:damage-after-field} 展开，可把关键中间量压缩成表~\ref{tab:damage-example-field}。

\begin{table}[H]
    \centering
    \caption{算例 1：属性克制与土场地的联合放大}
    \label{tab:damage-example-field}
    \begin{tabularx}{\textwidth}{>{\raggedright\arraybackslash}p{3.0cm} X >{\raggedleft\arraybackslash}p{2.7cm}}
        \toprule
        步骤 & 计算内容 & 结果 \\
        \midrule
        基础物伤 & 已知 $damage_0$ & $120$ \\
        属性克制主体 & $\bm{p}^{\mathsf T}M\bm{q}$ & $12000$ \\
        属性克制后伤害 & $120\times 12000/10000$ & $144$ \\
        攻方土场倍率 & $0.5+0.5\times 60\times 80/10000$ & $0.74$ \\
        守方土场倍率 & $0.5+0.5\times 10\times 80/10000$ & $0.54$ \\
        场地修正后伤害 & $144\times 0.74/0.54$ & $\approx 197.33$ \\
        \bottomrule
    \end{tabularx}
\end{table}

这个例子说明，属性克制已经先把伤害从 $120$ 推高到 $144$，而土场地又因为攻击方土属性显著高于防守方土属性，继续把伤害推到约 $197.33$。也就是说，场地收益并不是单独存在的，它会在属性克制已经放大的基础上继续按攻守比值放大。

\paragraph{例 2：属性反转同时改变场地收益与结界减伤}
设基础物伤 $damage_0=100$，防守方为纯无属性 $\bm{q}=(0,0,0,0,100)$，当前为火场地、场地强度 $s=100$，且防守方身上存在土结界。攻击方反转前的当前固定属性为 $\bm{p}=(80,20,10,5,0)$，按第~\ref{chap:dynamic-attr} 章的反转映射，反转后的当前固定属性变为 $\bm{p}^{(rev)}=(10,5,80,20,0)$。关键中间量汇总见表~\ref{tab:damage-example-reverse}。

\begin{table}[H]
    \centering
    \caption{算例 2：属性反转如何同时改写场地与结界}
    \label{tab:damage-example-reverse}
    \begin{tabularx}{\textwidth}{>{\raggedright\arraybackslash}p{3.1cm} >{\raggedleft\arraybackslash}X >{\raggedleft\arraybackslash}X}
        \toprule
        步骤 & 反转前 & 反转后 \\
        \midrule
        攻方当前属性 & $(80,20,10,5,0)$ & $(10,5,80,20,0)$ \\
        属性克制后伤害 $damage_{\mathrm{attr}}$ & $172.5$ & $172.5$ \\
        火场攻方倍率 & $0.55$ & $0.9$ \\
        火场守方倍率 & $0.5$ & $0.5$ \\
        场地修正后伤害 $damage_1$ & $189.75$ & $310.5$ \\
        土结界读取的攻方土属性 & $80$ & $10$ \\
        结界修正后伤害 $damage_2$ & $113.85$ & $294.975$ \\
        \bottomrule
    \end{tabularx}
\end{table}

这个例子非常典型：反转前后，属性克制主体本身没有变，但火场地倍率和土结界减伤同时被改写，最终伤害仍然从 $113.85$ 大幅跳到约 $294.98$。因此，实战里看到“反转后伤害忽然暴涨”，并不一定意味着攻防差或会心出了问题，很多时候只是属性向量在场地与结界两个层面同时改了方向。

\subsection{把属性层并入最终伤害顺序}

把本章前文与本节合并后，当前主分支下的物理伤害顺序可以整理为
\begin{equation}
    damage_0
    \xrightarrow{\text{属性克制矩阵}}
    damage_{\mathrm{attr}}
    \xrightarrow{\text{场地比值}}
    damage_1
    \xrightarrow{\text{四属性结界}}
    damage_2
    \xrightarrow{\text{额外扰动}}
    damage_3.
\end{equation}
其中若编译启用了额外伤害与额外减伤系统，则附加扰动写成
\begin{equation}
    \Delta_{extra} = \operatorname{RAND}(0.3P_a,P_a)-\operatorname{RAND}(0.3P_d,P_d),
\end{equation}
并有
\begin{equation}
    damage_3 = \max\{damage_2 + \Delta_{extra},0\}.
\end{equation}

这样一来，完整的物伤主体就不应再被理解成“攻防差后乘一个属性修正”。更准确地说，它是“\textbf{攻防分段函数 + 五维属性双线性克制 + 场地攻守比值 + 结界线性减伤 + 额外扰动}”的串联结构。特别是在骑宠环境下，攻击、防御会跟随骑宠混合，但属性向量仍读取当前单位自身工作值；这意味着“骑宠面板变了”与“属性克制结果变了”在源码里是两件相互独立的事情。

\section{会心额外伤害}

会心成立之后，伤害并不是直接乘以某个固定倍率。对当前实现而言，若武器不是弓，则系统先按普通流程算出基础物伤，再额外增加一段由防守方防御与双方等级共同决定的附加伤害。记普通物伤结果为 $damage$，攻击方等级与防守方等级分别为 $L_a$ 与 $L_d$，则会心附加后的结果可整理为
\begin{equation}
    damage_{crit} = damage + 0.5\,DEF_c\frac{L_a}{L_d},
    \label{eq:crit-damage}
\end{equation}
其中 $DEF_c$ 表示防守方的防御力工作值。式~\eqref{eq:crit-damage} 说明，会心本质上是“普通物伤 + 一段额外防御比例项”，而不是传统 RPG 中常见的“普通物伤乘一个暴击倍数”。

这一点对实战理解非常关键。若玩家只根据最终伤害浮动去猜测机制，往往会误以为会心是一个乘法型收益；但在当前源码口径下，会心的附加量与防守方防御及双方等级比直接相关，因此对不同目标的收益会显著不同。防高、等级低的目标可能吃到更大的会心附加，而防低、等级接近的目标则未必体现出同样夸张的增幅。进一步地，弓虽然仍可判出会心事件，但不会调用这段附加伤害函数，因此弓系配置在概率层面的会心和在爆发层面的会心并不是同一件事。

\section{变量影响分析}

\subsection{攻击对伤害的边际收益是分段的}

由式~\eqref{eq:damage-stage1}--\eqref{eq:damage-stage3} 可见，攻击力提升的边际收益并不是全局平滑递增的。当 $A<D$ 时，提高攻击的主要效果只是让随机区间从“几乎固定为 $0$ 或 $1$”逐步靠近第二段，边际收益非常弱；当 $D\le A < 8D/7$ 时，伤害上界随 $A/16$ 线性增加，但仍停留在小区间随机；只有跨过 $8D/7$ 这条阈值后，攻击提升才开始几乎按 $2(A-D)$ 的斜率进入主导区间。因此，若某个配置长期处在第二段附近，那么继续加攻击的体感收益可能远低于玩家的线性预期。

\subsection{骑宠远程与骑宠近战的收益结构不同}

由式~\eqref{eq:ride-atk-ranged} 与式~\eqref{eq:ride-atk-melee} 可知，骑宠远程与骑宠近战并不是“同一公式换个武器名字”。远程侧更偏人物攻击，近战侧则让人物和宠物都以较高权重进入攻击项；同时，敏捷侧又分别服从式~\eqref{eq:ride-quick-ranged} 与式~\eqref{eq:ride-quick-melee}。这意味着在骑宠配置下，“提高宠物攻击”与“提高宠物敏捷”的收益取决于你最终打算走远程还是近战，不能脱离武器类型单独讨论。

\subsection{防御与石化、忽视防御等状态的作用顺序}

防御量在进入伤害三段式之前，还可能受到铁壁防御、石化、无装备防御和忽视防御等状态影响。由于这些修正发生在不同位置，它们的实际效果并不相同。石化会直接把防御翻倍，从而整体抬高第三段门槛；忽视防御则在防御量已经形成之后再按比例削减，因此更接近对门槛的整体下压；若存在无装备防御类技能，则防御甚至会被直接替换为另一套工作值。由此可见，战斗中的“防御变化”并不只是某个数值加减，而是可能改变整条分段函数所处区间。

综上，骑宠与物伤系统共同呈现出一个高度离散、分段且带路径依赖的结构。攻击、防御、武器类型、骑宠状态、属性结界和会心并不是独立叠加，而是按特定顺序进入不同层级的公式。对于论文后续章节而言，这一结构恰好提供了一个桥梁：前文宠物成长与喂药所产生的数值差异，最终正是通过本章这些分段函数与混合公式，转化为实战中的先手、压制、爆发与承伤差异。
